\documentclass[12pt]{ximera}
\typeout{Start loading xmPreamble.tex}%

% Add here extra macro's that are loaded automatically by all documents of claas 'ximera' or 'xourse' in this repo

%%
%%  Example:
%%
% \newcommand{\R}{\mathbb{R}}
\usepackage{tikz}
\usetikzlibrary{positioning, arrows.meta}
\usepackage{multicol}
\usepackage{enumitem}
\usepackage{caption}
\usepackage{amsmath}
\usepackage{amsthm}
\usepackage{fontenc}

%% ---------------------------------------------------------------
%% Shared worksheet macros.
%%
%% These came from the Summer 2026 worksheet set, where each worksheet carried its
%% own copy. They have to live here instead: a xourse inlines every activity into a
%% single document, so duplicate \newcommand definitions across activities collide,
%% and the xourse gobbles \input so a per-topic macro file never loads.
%%
%% They use \providecommand, NOT \newcommand. ximera.cls loads this file automatically,
%% and every activity also carries an explicit \typeout{Start loading xmPreamble.tex}%

% Add here extra macro's that are loaded automatically by all documents of claas 'ximera' or 'xourse' in this repo

%%
%%  Example:
%%
% \newcommand{\R}{\mathbb{R}}
\usepackage{tikz}
\usetikzlibrary{positioning, arrows.meta}
\usepackage{multicol}
\usepackage{enumitem}
\usepackage{caption}
\usepackage{amsmath}
\usepackage{amsthm}
\usepackage{fontenc}

%% ---------------------------------------------------------------
%% Shared worksheet macros.
%%
%% These came from the Summer 2026 worksheet set, where each worksheet carried its
%% own copy. They have to live here instead: a xourse inlines every activity into a
%% single document, so duplicate \newcommand definitions across activities collide,
%% and the xourse gobbles \input so a per-topic macro file never loads.
%%
%% They use \providecommand, NOT \newcommand. ximera.cls loads this file automatically,
%% and every activity also carries an explicit \typeout{Start loading xmPreamble.tex}%

% Add here extra macro's that are loaded automatically by all documents of claas 'ximera' or 'xourse' in this repo

%%
%%  Example:
%%
% \newcommand{\R}{\mathbb{R}}
\usepackage{tikz}
\usetikzlibrary{positioning, arrows.meta}
\usepackage{multicol}
\usepackage{enumitem}
\usepackage{caption}
\usepackage{amsmath}
\usepackage{amsthm}
\usepackage{fontenc}

%% ---------------------------------------------------------------
%% Shared worksheet macros.
%%
%% These came from the Summer 2026 worksheet set, where each worksheet carried its
%% own copy. They have to live here instead: a xourse inlines every activity into a
%% single document, so duplicate \newcommand definitions across activities collide,
%% and the xourse gobbles \input so a per-topic macro file never loads.
%%
%% They use \providecommand, NOT \newcommand. ximera.cls loads this file automatically,
%% and every activity also carries an explicit \input{../xmPreamble}, so this file is
%% read twice. That was harmless while it held only \usepackage lines (LaTeX skips a
%% package it has already loaded) but a second \newcommand is a hard error.
%%
%%   \blank[width]        fill-in-the-blank rule (default 2.5cm)
%%   \showwork            "show and explain your work" reminder
%%   \showworkline        ... plus which schedule line was used
%%   \showworkyear        ... plus which year's schedule was used
%%   \schedhead{...}      centred bold caption above a tiered-rate table
%%   \samesquares[scale]{left}{right}   two equal-area squares, labelled
%%   \samecubes[scale]{left}{right}     two equal-volume cubes, labelled
%% ---------------------------------------------------------------

\providecommand{\blank}[1][2.5cm]{\underline{\hspace{#1}}}
\providecommand{\showwork}{\textbf{REMEMBER TO SHOW AND EXPLAIN YOUR WORK.}}
\providecommand{\showworkline}{\textbf{REMEMBER TO SHOW AND EXPLAIN YOUR WORK.
  Explanation should include how and why you know which line of the
  schedule was used to calculate this amount.}}
\providecommand{\showworkyear}{\begin{sloppypar}\textbf{REMEMBER TO SHOW AND
  EXPLAIN YOUR WORK. Explanation should include how and why you know which
  line of the year's tax schedule was used to calculate this tax
  amount.}\end{sloppypar}}
\providecommand{\schedhead}[1]{\begin{center}\textbf{#1}\end{center}}

\providecommand{\samesquares}[3][1]{%
\begin{center}
{\footnotesize These squares are the \textbf{same size}.}

\vspace{1.5mm}
\begin{tikzpicture}[scale=#1,every node/.style={scale=#1}]
  \draw (0,0) rectangle (2.4,2.4);
  \node[anchor=north west] at (0.15,2.25) {Area:};
  \node[anchor=east]  at (-0.15,1.2) {#2};
  \node[anchor=north] at (1.2,-0.15) {#2};
  \begin{scope}[xshift=4.6cm]
    \draw (0,0) rectangle (2.4,2.4);
    \node[anchor=north west] at (0.15,2.25) {Area:};
    \node[anchor=east]  at (-0.15,1.2) {#3};
    \node[anchor=north] at (1.2,-0.15) {#3};
  \end{scope}
\end{tikzpicture}
\end{center}}

\providecommand{\samecubes}[3][1]{%
\begin{center}
{\footnotesize These cubes are the \textbf{same size}.}

\vspace{1.5mm}
\begin{tikzpicture}[scale=#1,every node/.style={scale=#1}]
  \foreach \dx in {0,5.2} {
    \begin{scope}[xshift=\dx cm]
      \draw (0,0) rectangle (2.0,2.0);
      \draw (0,2.0) -- (0.65,2.6) -- (2.65,2.6) -- (2.0,2.0);
      \fill[black!12] (2.0,0) -- (2.65,0.6) -- (2.65,2.6) -- (2.0,2.0) -- cycle;
      \draw (2.0,0) -- (2.65,0.6) -- (2.65,2.6) -- (2.0,2.0);
      \node[anchor=north west] at (0.1,1.9) {Volume:};
    \end{scope}
  }
  \node[anchor=east]  at (-0.15,1.0) {#2};
  \node[anchor=north] at (1.0,-0.15) {#2};
  \node[anchor=west]  at (2.25,0.4) {#2};
  \node[anchor=east]  at (5.05,1.0) {#3};
  \node[anchor=north] at (6.2,-0.15) {#3};
  \node[anchor=west]  at (7.45,0.4) {#3};
\end{tikzpicture}
\end{center}}

%% NOTE: do NOT add \usepackage to individual activity .tex files.
%% When a xourse inlines an activity it gobbles \documentclass and \input, but a bare
%% \usepackage still executes -- after \begin{document} -- and errors out.
%% All shared packages and macros belong here.
%%
%% ximera.cls already loads: amsmath, amssymb, amsthm, cancel, enumitem, graphicx,
%% hyperref, listings, pgfplots, tikz, url, xcolor. No need to re-load those.
, so this file is
%% read twice. That was harmless while it held only \usepackage lines (LaTeX skips a
%% package it has already loaded) but a second \newcommand is a hard error.
%%
%%   \blank[width]        fill-in-the-blank rule (default 2.5cm)
%%   \showwork            "show and explain your work" reminder
%%   \showworkline        ... plus which schedule line was used
%%   \showworkyear        ... plus which year's schedule was used
%%   \schedhead{...}      centred bold caption above a tiered-rate table
%%   \samesquares[scale]{left}{right}   two equal-area squares, labelled
%%   \samecubes[scale]{left}{right}     two equal-volume cubes, labelled
%% ---------------------------------------------------------------

\providecommand{\blank}[1][2.5cm]{\underline{\hspace{#1}}}
\providecommand{\showwork}{\textbf{REMEMBER TO SHOW AND EXPLAIN YOUR WORK.}}
\providecommand{\showworkline}{\textbf{REMEMBER TO SHOW AND EXPLAIN YOUR WORK.
  Explanation should include how and why you know which line of the
  schedule was used to calculate this amount.}}
\providecommand{\showworkyear}{\begin{sloppypar}\textbf{REMEMBER TO SHOW AND
  EXPLAIN YOUR WORK. Explanation should include how and why you know which
  line of the year's tax schedule was used to calculate this tax
  amount.}\end{sloppypar}}
\providecommand{\schedhead}[1]{\begin{center}\textbf{#1}\end{center}}

\providecommand{\samesquares}[3][1]{%
\begin{center}
{\footnotesize These squares are the \textbf{same size}.}

\vspace{1.5mm}
\begin{tikzpicture}[scale=#1,every node/.style={scale=#1}]
  \draw (0,0) rectangle (2.4,2.4);
  \node[anchor=north west] at (0.15,2.25) {Area:};
  \node[anchor=east]  at (-0.15,1.2) {#2};
  \node[anchor=north] at (1.2,-0.15) {#2};
  \begin{scope}[xshift=4.6cm]
    \draw (0,0) rectangle (2.4,2.4);
    \node[anchor=north west] at (0.15,2.25) {Area:};
    \node[anchor=east]  at (-0.15,1.2) {#3};
    \node[anchor=north] at (1.2,-0.15) {#3};
  \end{scope}
\end{tikzpicture}
\end{center}}

\providecommand{\samecubes}[3][1]{%
\begin{center}
{\footnotesize These cubes are the \textbf{same size}.}

\vspace{1.5mm}
\begin{tikzpicture}[scale=#1,every node/.style={scale=#1}]
  \foreach \dx in {0,5.2} {
    \begin{scope}[xshift=\dx cm]
      \draw (0,0) rectangle (2.0,2.0);
      \draw (0,2.0) -- (0.65,2.6) -- (2.65,2.6) -- (2.0,2.0);
      \fill[black!12] (2.0,0) -- (2.65,0.6) -- (2.65,2.6) -- (2.0,2.0) -- cycle;
      \draw (2.0,0) -- (2.65,0.6) -- (2.65,2.6) -- (2.0,2.0);
      \node[anchor=north west] at (0.1,1.9) {Volume:};
    \end{scope}
  }
  \node[anchor=east]  at (-0.15,1.0) {#2};
  \node[anchor=north] at (1.0,-0.15) {#2};
  \node[anchor=west]  at (2.25,0.4) {#2};
  \node[anchor=east]  at (5.05,1.0) {#3};
  \node[anchor=north] at (6.2,-0.15) {#3};
  \node[anchor=west]  at (7.45,0.4) {#3};
\end{tikzpicture}
\end{center}}

%% NOTE: do NOT add \usepackage to individual activity .tex files.
%% When a xourse inlines an activity it gobbles \documentclass and \input, but a bare
%% \usepackage still executes -- after \begin{document} -- and errors out.
%% All shared packages and macros belong here.
%%
%% ximera.cls already loads: amsmath, amssymb, amsthm, cancel, enumitem, graphicx,
%% hyperref, listings, pgfplots, tikz, url, xcolor. No need to re-load those.
, so this file is
%% read twice. That was harmless while it held only \usepackage lines (LaTeX skips a
%% package it has already loaded) but a second \newcommand is a hard error.
%%
%%   \blank[width]        fill-in-the-blank rule (default 2.5cm)
%%   \showwork            "show and explain your work" reminder
%%   \showworkline        ... plus which schedule line was used
%%   \showworkyear        ... plus which year's schedule was used
%%   \schedhead{...}      centred bold caption above a tiered-rate table
%%   \samesquares[scale]{left}{right}   two equal-area squares, labelled
%%   \samecubes[scale]{left}{right}     two equal-volume cubes, labelled
%% ---------------------------------------------------------------

\providecommand{\blank}[1][2.5cm]{\underline{\hspace{#1}}}
\providecommand{\showwork}{\textbf{REMEMBER TO SHOW AND EXPLAIN YOUR WORK.}}
\providecommand{\showworkline}{\textbf{REMEMBER TO SHOW AND EXPLAIN YOUR WORK.
  Explanation should include how and why you know which line of the
  schedule was used to calculate this amount.}}
\providecommand{\showworkyear}{\begin{sloppypar}\textbf{REMEMBER TO SHOW AND
  EXPLAIN YOUR WORK. Explanation should include how and why you know which
  line of the year's tax schedule was used to calculate this tax
  amount.}\end{sloppypar}}
\providecommand{\schedhead}[1]{\begin{center}\textbf{#1}\end{center}}

\providecommand{\samesquares}[3][1]{%
\begin{center}
{\footnotesize These squares are the \textbf{same size}.}

\vspace{1.5mm}
\begin{tikzpicture}[scale=#1,every node/.style={scale=#1}]
  \draw (0,0) rectangle (2.4,2.4);
  \node[anchor=north west] at (0.15,2.25) {Area:};
  \node[anchor=east]  at (-0.15,1.2) {#2};
  \node[anchor=north] at (1.2,-0.15) {#2};
  \begin{scope}[xshift=4.6cm]
    \draw (0,0) rectangle (2.4,2.4);
    \node[anchor=north west] at (0.15,2.25) {Area:};
    \node[anchor=east]  at (-0.15,1.2) {#3};
    \node[anchor=north] at (1.2,-0.15) {#3};
  \end{scope}
\end{tikzpicture}
\end{center}}

\providecommand{\samecubes}[3][1]{%
\begin{center}
{\footnotesize These cubes are the \textbf{same size}.}

\vspace{1.5mm}
\begin{tikzpicture}[scale=#1,every node/.style={scale=#1}]
  \foreach \dx in {0,5.2} {
    \begin{scope}[xshift=\dx cm]
      \draw (0,0) rectangle (2.0,2.0);
      \draw (0,2.0) -- (0.65,2.6) -- (2.65,2.6) -- (2.0,2.0);
      \fill[black!12] (2.0,0) -- (2.65,0.6) -- (2.65,2.6) -- (2.0,2.0) -- cycle;
      \draw (2.0,0) -- (2.65,0.6) -- (2.65,2.6) -- (2.0,2.0);
      \node[anchor=north west] at (0.1,1.9) {Volume:};
    \end{scope}
  }
  \node[anchor=east]  at (-0.15,1.0) {#2};
  \node[anchor=north] at (1.0,-0.15) {#2};
  \node[anchor=west]  at (2.25,0.4) {#2};
  \node[anchor=east]  at (5.05,1.0) {#3};
  \node[anchor=north] at (6.2,-0.15) {#3};
  \node[anchor=west]  at (7.45,0.4) {#3};
\end{tikzpicture}
\end{center}}

%% NOTE: do NOT add \usepackage to individual activity .tex files.
%% When a xourse inlines an activity it gobbles \documentclass and \input, but a bare
%% \usepackage still executes -- after \begin{document} -- and errors out.
%% All shared packages and macros belong here.
%%
%% ximera.cls already loads: amsmath, amssymb, amsthm, cancel, enumitem, graphicx,
%% hyperref, listings, pgfplots, tikz, url, xcolor. No need to re-load those.

\title{In-Class: Income Taxes}
\begin{document}
\begin{abstract}
Using the 2019 tax schedule rather than the tax table: computing total tax, net
income, and effective tax rate; marginal income and marginal tax rate, including a
raise that crosses a bracket boundary; and working backwards from a tax paid to an
adjusted gross income.
\end{abstract}
\maketitle

\section*{Income Tax Activity}

In your Pre-Class Activity, you looked up how much needed to be paid in taxes from a
table of values. But where exactly do the values in the Tax Table come from? The
entire tax table can be described more concisely by a \emph{Tax Schedule} --- a table
of instructions for how to calculate taxes from the taxable income. Using these
instructions, we can see how the values in the table were determined.

The schedule is at the end of this activity; you may want to keep it open in another
window.

\subsection*{From income to tax}

\begin{example}
Using the tax schedule (not the tax table from the pre-class activity), calculate the
taxes that should be paid on an AGI of \$25,000.

\begin{center}
\begin{tabular}{|l|l|l|}
\hline
1040 Line(s) & Description & Amount \\
\hline
1 and 8b & Wages, salaries, and tips. Adjusted gross income (AGI) & \$25,000 \\
\hline
11a & Standard deduction and exemption (single, nondependent) & \$12,200 \\
\hline
11b & Subtract line 11a from line 8b. Taxable income & \$12,800 \\
\hline
16 & Total Tax (from schedule) & \$1,342 \\
\hline
 & Effective tax rate (tax as a percentage of all income) & $5.37\%$ \\
\hline
\end{tabular}
\end{center}

Taxable income: $\$25{,}000 - \$12{,}200 = \$12{,}800$.

Now look at the tax schedule. The taxable income is more than \$9,700 but less than
\$39,475, so we use the instruction ``Tax is \$970 plus $12\%$ of the amount over
\$9,700.''

\[
\begin{aligned}
\text{Taxable income over } \$9{,}700: &\quad \$12{,}800 - \$9{,}700 = \$3{,}100 \\
12\% \text{ of the amount over}: &\quad 0.12 \cdot \$3{,}100 = \$372 \\
\text{Total tax}: &\quad \$970 + \$372 = \$1{,}342
\end{aligned}
\]

Effective rate --- total tax is what percent of AGI? $\$1{,}342 = x \cdot \$25{,}000$,
so $x = \frac{\$1{,}342}{\$25{,}000} = 0.05368$, about $5.37\%$. This individual is
paying about $5.37\%$ of their AGI in taxes.
\end{example}

\begin{problem}
Now do the same for an AGI of \$65,000.

\begin{center}
\begin{tabular}{|l|l|l|}
\hline
1040 Line(s) & Description & Amount \\
\hline
1 and 8b & Wages, salaries, and tips. Adjusted gross income (AGI) & \$65,000 \\
\hline
11a & Standard deduction and exemption (single, nondependent) & \$12,200 \\
\hline
11b & Subtract line 11a from line 8b. Taxable income & \$$\answer{52800}$ \\
\hline
16 & Total Tax (from schedule) & \$$\answer[tolerance=0.5]{7474.50}$ \\
\hline
 & Effective tax rate & $\answer[tolerance=0.05]{11.50}$\% \\
\hline
\end{tabular}
\end{center}

\begin{solution}
Taxable income: $\$65{,}000 - \$12{,}200 = \$52{,}800$.

\$52,800 is more than \$39,475 but less than \$84,200, so we use ``Tax is \$4,543
plus $22\%$ of the amount over \$39,475.''
\[
\begin{aligned}
\$52{,}800 - \$39{,}475 &= \$13{,}325 \\
0.22 \cdot \$13{,}325 &= \$2{,}931.50 \\
\text{Total tax}: \quad \$4{,}543 + \$2{,}931.50 &= \$7{,}474.50
\end{aligned}
\]
Effective rate: $\frac{\$7{,}474.50}{\$65{,}000} = 0.1150$, about $11.50\%$.
\end{solution}
\end{problem}

\begin{problem}
From the Pre-Class Activity, record the taxes owed on \$25,000 and \$65,000 of income
according to the Tax \emph{Table}. How do those compare to what you just calculated
from the Tax \emph{Schedule}? Why is there a difference?

\begin{freeResponse}
\end{freeResponse}

The tax table gives a tax amount for every \$$\answer{50}$ increment of taxable income.

\begin{solution}
The table is built in fixed increments of taxable income and reports a single tax
figure for everyone inside each increment, so unless your taxable income lands exactly
where the increment does, the table value and the schedule value must differ slightly.
\end{solution}
\end{problem}

As you may have noticed, the table cannot be used for taxable incomes over \$100,000.
For incomes greater than that, other differences and rules may apply, but our main
focus is that since the table has no values above \$100,000, the Tax Schedule must be
used. We will ignore other possible complications by assuming the process is the same
even for these high incomes. (This is not a taxes class --- we are only after the
general mathematical ideas.)

\subsection*{Marginal income and marginal tax rate}

\begin{problem}
Use the tax schedule to compute the taxes and other table values beginning with a
salary of \$120,000, and again after a $3\%$ raise.

\begin{center}
\begin{tabular}{|l|l|l|}
\hline
Description & Before Raise & After Raise \\
\hline
Adjusted gross income (AGI) & \$120,000 & \$$\answer{123600}$ \\
\hline
Standard deduction and exemption & \$12,200 & \$12,200 \\
\hline
Taxable income & \$$\answer{107800}$ & \$$\answer{111400}$ \\
\hline
Total Tax (from schedule) & \$$\answer{20046}$ & \$$\answer{20910}$ \\
\hline
Net Income (after Fed. Inc. Tax) & \$$\answer{99954}$ & \$$\answer{102690}$ \\
\hline
Effective tax rate & $\answer[tolerance=0.05]{16.71}$\% & $\answer[tolerance=0.05]{16.92}$\% \\
\hline
Additional (marginal) income & \multicolumn{2}{c|}{\$$\answer{3600}$} \\
\hline
Additional (marginal) tax & \multicolumn{2}{c|}{\$$\answer{864}$} \\
\hline
Marginal tax rate & \multicolumn{2}{c|}{$\answer{24}$\%} \\
\hline
\end{tabular}
\end{center}

\begin{solution}
Income after the $3\%$ raise:
$\$120{,}000 + 0.03 \cdot \$120{,}000 = \$120{,}000 + \$3{,}600 = \$123{,}600$.

Taxable income: $\$120{,}000 - \$12{,}200 = \$107{,}800$ before, and
$\$123{,}600 - \$12{,}200 = \$111{,}400$ after.

Both \$107,800 and \$111,400 are more than \$84,200 but less than \$160,725, so both
use ``Tax is \$14,382 plus $24\%$ of the amount over \$84,200.''
\[
\begin{aligned}
\text{Before}: \quad \$107{,}800 - \$84{,}200 = \$23{,}600, &\quad 0.24 \cdot \$23{,}600 = \$5{,}664 \\
\text{After}: \quad \$111{,}400 - \$84{,}200 = \$27{,}200, &\quad 0.24 \cdot \$27{,}200 = \$6{,}528
\end{aligned}
\]
Total tax: $\$14{,}382 + \$5{,}664 = \$20{,}046$ before,
$\$14{,}382 + \$6{,}528 = \$20{,}910$ after.

Net income: $\$120{,}000 - \$20{,}046 = \$99{,}954$ and
$\$123{,}600 - \$20{,}910 = \$102{,}690$.

Effective rates: $\frac{\$20{,}046}{\$120{,}000} = 0.1671$ and
$\frac{\$20{,}910}{\$123{,}600} = 0.1692$.

Marginal income: $\$123{,}600 - \$120{,}000 = \$3{,}600$.
Marginal tax: $\$20{,}910 - \$20{,}046 = \$864$.
Marginal tax rate: $\frac{\$864}{\$3{,}600} = 0.24$, so $24\%$ --- the extra income
from the raise is taxed at $24\%$.
\end{solution}
\end{problem}

\begin{problem}
Look back at the marginal tax rates for the various raises and record them:

$3\%$ raise on \$25,000 income: marginal rate $\answer{12}$\%

$3\%$ raise on \$65,000 income: marginal rate $\answer{22}$\%

$3\%$ raise on \$120,000 income: marginal rate $\answer{24}$\%

Where do you see these marginal tax values in the tax schedule? Explain why these
values should be the same.

\begin{hint}
Look at the taxable incomes before and after each raise and identify which bracket
each falls into. For the \$25,000 income: which bracket is the taxable income in
before the raise? Which after?
\end{hint}

\begin{freeResponse}
\end{freeResponse}

\begin{solution}
In each case the raise is small enough that the taxable income stays inside the
\emph{same} bracket before and after. Every extra dollar is therefore taxed at that
bracket's stated rate, so the marginal tax rate is exactly the bracket percentage
from the schedule.
\end{solution}
\end{problem}

\begin{problem}
Now let's explore what happens if a raise \emph{does} move income into the next
bracket. Use the tax schedule beginning with a salary of \$96,000 and a $3\%$ raise.

\begin{center}
\begin{tabular}{|l|l|l|}
\hline
Description & Before Raise & After Raise \\
\hline
Adjusted gross income (AGI) & \$96,000 & \$$\answer{98880}$ \\
\hline
Standard deduction and exemption & \$12,200 & \$12,200 \\
\hline
Taxable income & \$$\answer{83800}$ & \$$\answer{86680}$ \\
\hline
Total Tax (from schedule) & \$$\answer[tolerance=0.5]{14294.50}$ & \$$\answer[tolerance=0.5]{14977.20}$ \\
\hline
Net Income (after Fed. Inc. Tax) & \$$\answer[tolerance=0.5]{81705.50}$ & \$$\answer[tolerance=0.5]{83902.80}$ \\
\hline
Effective tax rate & $\answer[tolerance=0.05]{14.89}$\% & $\answer[tolerance=0.05]{15.15}$\% \\
\hline
Additional (marginal) income & \multicolumn{2}{c|}{\$$\answer{2880}$} \\
\hline
Additional (marginal) tax & \multicolumn{2}{c|}{\$$\answer[tolerance=0.5]{682.70}$} \\
\hline
Marginal tax rate & \multicolumn{2}{c|}{$\answer[tolerance=0.05]{23.70}$\%} \\
\hline
\end{tabular}
\end{center}

\begin{solution}
Income after the $3\%$ raise:
$\$96{,}000 + 0.03 \cdot \$96{,}000 = \$96{,}000 + \$2{,}880 = \$98{,}880$.

Taxable income: $\$96{,}000 - \$12{,}200 = \$83{,}800$ before,
$\$98{,}880 - \$12{,}200 = \$86{,}680$ after.

Before the raise, \$83,800 is more than \$39,475 but less than \$84,200, so we use
``Tax is \$4,543 plus $22\%$ of the amount over \$39,475.'' After the raise,
\$86,680 is more than \$84,200 but less than \$160,725, so we use ``Tax is \$14,382
plus $24\%$ of the amount over \$84,200.'' \emph{The bracket changed.}

\begin{center}
\begin{tabular}{|l|l|l|}
\hline
 & Before Raise & After Raise \\
\hline
Amount over & $\$83{,}800 - \$39{,}475 = \$44{,}325$ & $\$86{,}680 - \$84{,}200 = \$2{,}480$ \\
\hline
Tax on amount over & $0.22 \cdot \$44{,}325 = \$9{,}751.50$ & $0.24 \cdot \$2{,}480 = \$595.20$ \\
\hline
Total tax & $\$4{,}543 + \$9{,}751.50 = \$14{,}294.50$ & $\$14{,}382 + \$595.20 = \$14{,}977.20$ \\
\hline
Net income & $\$96{,}000 - \$14{,}294.50 = \$81{,}705.50$ & $\$98{,}880 - \$14{,}977.20 = \$83{,}902.80$ \\
\hline
Effective rate & $\frac{\$14{,}294.50}{\$96{,}000} = 14.89\%$ & $\frac{\$14{,}977.20}{\$98{,}880} = 15.15\%$ \\
\hline
\end{tabular}
\end{center}

Marginal income: $\$98{,}880 - \$96{,}000 = \$2{,}880$.
Marginal tax: $\$14{,}977.20 - \$14{,}294.50 = \$682.70$.
Marginal tax rate: $\frac{\$682.70}{\$2{,}880} = 0.2370$, so $23.70\%$.
\end{solution}
\end{problem}

\begin{problem}
This marginal tax rate does not match any single value in the schedule, as it did
before. What two values does it fall between? Explain why that makes sense.

It falls between $\answer{22}$\% and $\answer{24}$\%.

\begin{freeResponse}
\end{freeResponse}

\begin{solution}
The before-raise and after-raise incomes did not fall in the same bracket. Before the
raise we were in the $22\%$ bracket; after, in the $24\%$ bracket. Since part of the
raise is taxed at $22\%$ and part at $24\%$, the overall marginal rate lands between
the two.
\end{solution}
\end{problem}

\subsection*{Working backwards}

So far we have always calculated the taxes from the income. But we can work in the
other direction: from how much someone paid in taxes we can get an idea of what their
income might have been. We do have to make some assumptions --- filing single, no
earned income credit, and so on --- the same assumptions we started with in the
Pre-Class Activity.

\begin{problem}
Suppose someone paid \$6,700 in taxes. What could this person's adjusted gross income
be? Try looking at your process for calculating taxes above, and think about
``undoing'' it step by step.

\begin{center}
\begin{tabular}{|l|l|}
\hline
Description & Amount \\
\hline
Adjusted gross income (AGI) & \$$\answer[tolerance=1]{61479.55}$ \\
\hline
Standard deduction and exemption & \$12,200 \\
\hline
Taxable income & \$$\answer[tolerance=1]{49279.55}$ \\
\hline
Total Tax (from schedule) & \$6,700 \\
\hline
Effective tax rate & $\answer[tolerance=0.05]{10.90}$\% \\
\hline
\end{tabular}
\end{center}

\begin{hint}
Which schedule line could produce \$6,700? It must be the line whose flat amount is
below \$6,700 but whose \emph{next} line's flat amount is above it.
\end{hint}

\begin{solution}
\$6,700 must come from the instruction ``\$4,543 plus $22\%$ of the amount over
\$39,475,'' because \$6,700 is more than the flat rate \$4,543 but less than the flat
rate \$14,382 of the next line.

So $\$4{,}543 + 0.22 \cdot (\text{taxable income} - \$39{,}475) = \$6{,}700$.
Subtracting \$4,543 from both sides,
\[
0.22 (\text{taxable income} - \$39{,}475) = \$6{,}700 - \$4{,}543 = \$2{,}157,
\]
and dividing by $0.22$,
\[
\text{taxable income} - \$39{,}475 = \frac{\$2{,}157}{0.22} = \$9{,}804.55.
\]
So taxable income $= \$9{,}804.55 + \$39{,}475 = \$49{,}279.55$, and
\[
\text{AGI} = \text{taxable income} + \text{standard deduction}
= \$49{,}279.55 + \$12{,}200 = \$61{,}479.55,
\]
about \$61,480.

Effective rate: $\$6{,}700 = x \cdot \$61{,}480$, so
$x = \frac{\$6{,}700}{\$61{,}480} = 0.1090$, about $10.90\%$ of their AGI.
\end{solution}
\end{problem}

\begin{problem}
What may help for these problems is to think about how you would calculate the taxes
if you knew the taxable income. Let $x$ be the taxable income, and suppose $x$ is
between \$9,700 and \$39,475.

\begin{prompt}
\textbf{(a)} What must the adjusted gross income be, in terms of $x$?
\begin{freeResponse}
\end{freeResponse}

\begin{solution}
$\text{AGI} = x + \$12{,}200$, so the AGI is between $\$21{,}900$ and $\$51{,}675$.
\end{solution}
\end{prompt}

\begin{prompt}
\textbf{(b)} Use $x$ to write a function for the amount of taxes to be paid on $x$.
\begin{freeResponse}
\end{freeResponse}

\begin{solution}
$T(x) = 970 + 0.12(x - 9700)$ for $9700 \le x \le 39475$. Expanding,
$T(x) = 0.12x - 194$ --- a linear function, which is why each bracket is a straight
line segment.
\end{solution}
\end{prompt}
\end{problem}

\subsection*{Single Taxable Income Tax Brackets and Rates, 2019}

\begin{center}
\begin{tabular}{|l|l|l|}
\hline
If taxable income is over & but not over & the tax is: \\
\hline
\$0 & \$9,700 & $10\%$ of the amount over \$0 \\
\hline
\$9,700 & \$39,475 & \$970 plus $12\%$ of the amount over \$9,700 \\
\hline
\$39,475 & \$84,200 & \$4,543 plus $22\%$ of the amount over \$39,475 \\
\hline
\$84,200 & \$160,725 & \$14,382 plus $24\%$ of the amount over \$84,200 \\
\hline
\$160,725 & \$204,100 & \$32,748 plus $32\%$ of the amount over \$160,725 \\
\hline
\$204,100 & \$510,300 & \$46,628 plus $35\%$ of the amount over \$204,100 \\
\hline
\$510,300 & no limit & \$153,798 plus $37\%$ of the amount over \$510,300 \\
\hline
\end{tabular}
\end{center}

\end{document}
